%----------------------------------------------------------------------------------------
% OPPERVLAKKEN EN RUIMTEKROMMEN. Alle opgaven zijn eigen werk.
%
% Opbouw: Probleem, Start, leerstof, en op een nieuwe pagina de Opdrachten, Koppel en Test jezelf.
% Labels van dit hoofdstuk beginnen met h2-. De berekende figuren (PNG) komen uit afbeeldingen/maak-afbeeldingen.py.
%----------------------------------------------------------------------------------------

\mychapter{2}{Oppervlakken en ruimtekrommen}

\begin{probleem}
\vspace{-1mm}\naast{Een wenteltrap draait rond een zuil met een straal van 0,2~m. De treden zijn 1,2~m breed, dus de buitenrand ligt op 1,4~m van de as. Bij één volle draai stijgt de trap 3~m. Langs de buitenrand loopt een leuning. Hoe lang is de leuning voor één volle draai? En waarom is zo'n trap aan de binnenkant veel steiler dan aan de buitenkant: hoe steil is hij aan beide kanten, in graden?}{\begin{tikzpicture}[zicht=.85,font=\scriptsize]
  \draw[lijngrijs] (0,-.2,0) -- (0,-.2,4.2) (0,.2,0) -- (0,.2,4.2);
  \foreach \t in {0,30,...,420}{\draw[grijs] ({.2*cos(\t)},{.2*sin(\t)},{3*\t/360}) -- ({1.4*cos(\t)},{1.4*sin(\t)},{3*\t/360});}
  \draw[accent,line width=1.4pt] plot[domain=0:420,samples=90,smooth,variable=\t] ({1.4*cos(\t)},{1.4*sin(\t)},{3*\t/360+.45});
  \draw[accent!60] plot[domain=0:420,samples=90,smooth,variable=\t] ({1.4*cos(\t)},{1.4*sin(\t)},{3*\t/360});
\end{tikzpicture}}
\end{probleem}

\lkr{Leerplan Complementair: LPD 12, met de wenken over ICT, de helicoïde, de torus en de sombrero, en over de schroeflijn (helix) en de kubische schroeflijn. Voorkennis uit pakket 4: Ruimtemeetkunde (parametervergelijkingen van een rechte en een vlak, een rechte als snijlijn van twee vlakken, vlakken evenwijdig met een as), goniometrie ($\cos^2 t + \sin^2 t = 1$) en de parametervergelijking van een cirkel en een ellips (Kegelsneden, hoofdstuk 3), en hoofdstuk 1. Wie de module Krommen en meetkundige plaatsen volgde, kent parametervoorstellingen van vlakke krommen, de snelheidsvector en poolcoördinaten: dit hoofdstuk breidt dat uit naar de ruimte, maar steunt er niet op. ICT: GeoGebra 3D tekent elk oppervlak van dit hoofdstuk (\textit{Oppervlak} en \textit{Kromme}); in de les volstaan de figuren in het boek. Extra's: de kromme van Viviani, uitrollen van een cilinder, en de stelling dat een vlak de kubische schroeflijn hoogstens drie keer snijdt.}

In hoofdstuk 1 had de bol één vergelijking, en daarmee kon je alles berekenen. Maar hoe beschrijf je de vorm van een wenteltrap, een donut of een golvend dak? En hoe beschrijf je een kromme die niet in een vlak ligt, zoals de leuning van een wenteltrap? In dit hoofdstuk leer je twee manieren: een cartesiaanse vergelijking, die zegt of een punt op het oppervlak ligt, en een parametervoorstelling, die alle punten één voor één oplevert. Een computer tekent oppervlakken met de tweede.

\begin{oefeningen}[Start]
\oef{h2-start-1}

\oef{h2-verken}
\end{oefeningen}

%========================================================================================
\section{Oppervlakken met een vergelijking}

\subsection{Doorsneden met vlakken}

Een vlak heeft een vergelijking van de eerste graad, een bol een vergelijking van de tweede graad. Elke vergelijking in $x$, $y$ en $z$ beschrijft een verzameling punten in de ruimte, en meestal is dat een oppervlak. Maar hoe zie je aan $z = x^2 + y^2$ welke vorm het heeft? Op een landkaart toont een tekenaar een berglandschap met hoogtelijnen: de punten die even hoog liggen. Zo snijd je een oppervlak met horizontale vlakken, en kijk je welke krommen je krijgt.

\begin{theorie}
Een \textbf{\textcolor{accent}{cartesiaanse vergelijking}} van een oppervlak is een vergelijking $\formule{F(x, y, z) = 0}$ waaraan de coördinaten van de punten van het oppervlak voldoen, en alleen die.

De doorsnede met een horizontaal vlak $\formule{z = c}$ heet een \textbf{\textcolor{accent}{niveaukromme}}. Je vindt ze door $z = c$ in de vergelijking in te vullen.
\end{theorie}

Ook verticale doorsneden helpen: vul $x = 0$ in, of $y = 0$, en je ziet het profiel van het oppervlak in een coördinaatvlak. Een paar doorsneden samen geven een goed beeld.

\voorbeeld{}{Beschrijf het oppervlak $z = x^2 + y^2$ met doorsneden.\label{vb:h2-doorsneden}}{%
\uitleg{We snijden met horizontale vlakken $z = c$, en met het $yz$-vlak $x = 0$.}
\naast{\begin{stappenvrij}{4.4cm}{4.9cm}
c < 0 & $x^2 + y^2 = c$ heeft geen oplossing: het oppervlak ligt boven het $xy$-vlak. \\[1ex]
c = 0 & Alleen de oorsprong. \\[1ex]
c > 0\colon\ x^2 + y^2 = c & Een cirkel met straal $\sqrt c$ op hoogte $c$. \\[1ex]
x = 0\colon\ z = y^2 & Het profiel is een parabool.
\end{stappenvrij}}{\begin{tikzpicture}[zicht=.6,font=\scriptsize]
  \beeld{kom.png}{-1.010}{-.068}{2.020}
  \draw[->,lijngrijs] (0,0,2.56) -- (0,0,3.1) node[above,text=black] {$z$};
\end{tikzpicture}\quad\begin{tikzpicture}[scale=.55,font=\scriptsize]
  \draw[->,lijngrijs] (-1.8,0) -- (1.9,0) node[right,text=black] {$x$};
  \draw[->,lijngrijs] (0,-1.8) -- (0,1.9) node[above,text=black] {$y$};
  \foreach \c in {.25,.5,...,2.5}{\draw[accent] (0,0) circle ({sqrt(\c)});}
\end{tikzpicture}}
De niveaukrommen zijn cirkels die steeds verder uit elkaar liggen (rechts, voor $c = 0{,}25$; $0{,}5$; $\ldots$; $2{,}5$): het oppervlak wordt steeds steiler. Het is een kom, met in elk verticaal vlak door de $z$-as een parabool. In hoofdstuk 3 heet het een \textit{paraboloïde}.}

\begin{werkwijze}
Een oppervlak met een vergelijking verkennen: vul $z = c$ in voor een paar waarden van $c$ en herken de niveaukrommen. Vul $x = 0$ en $y = 0$ in voor de profielen. Teken ze, en zet ze samen in een ruimtelijke schets.
\end{werkwijze}

\nujij{Welke niveaukrommen heeft het oppervlak $z = 2x + y$? Welk oppervlak is het?}{De rechten $2x + y = c$, evenwijdig: het is een vlak.}

\begin{oefeningen}
\oef{h2-doorsneden}
\end{oefeningen}

\subsection{Cilinders}

Wat stelt $x^2 + y^2 = 4$ voor? In het vlak is het een cirkel. Maar in de ruimte staat er geen $z$ in de vergelijking, en dus legt ze $z$ niet vast. Ligt $(2, 0, 0)$ op het oppervlak, dan ook $(2, 0, 5)$ en $(2, 0, -3)$: de hele verticale rechte door dat punt. Het oppervlak bestaat uit verticale rechten door de punten van de cirkel. Dat is een cilinder.

\begin{theorie}
\begin{minipage}[t]{.68\linewidth}
\vspace{0pt}
Een \textbf{\textcolor{accent}{cilinder}} is een oppervlak dat bestaat uit evenwijdige rechten, de \textbf{\textcolor{accent}{beschrijvenden}}, door de punten van een kromme, de \textbf{\textcolor{accent}{richtkromme}}.

\begin{stelling}{Vergelijking zonder z}
Een vergelijking $F(x, y) = 0$ waarin $z$ niet voorkomt, stelt in de ruimte een cilinder voor met verticale beschrijvenden. De richtkromme is de kromme $F(x, y) = 0$ in het $xy$-vlak.
\end{stelling}
\end{minipage}\hfill
\begin{minipage}[t]{.28\linewidth}
\vspace{0pt}
\centering
\begin{tikzpicture}[zicht=.55,font=\scriptsize]
  \foreach \t in {-60,-30,...,120}{\draw[accent!50] ({2*cos(\t)},{2*sin(\t)},0) -- ({2*cos(\t)},{2*sin(\t)},3);}
  \cirkelzicht{.55}{0}{0}{0}{2}
  \cirkelzicht[thick]{.55}{0}{0}{3}{2}
  \draw[accent,thick] (2,0,0) -- (2,0,3);
  \draw[accent,thick] ([xshift=-1.1cm]0,0,0) -- ([xshift=-1.1cm]0,0,3);
  \draw[accent,thick] ([xshift=1.1cm]0,0,0) -- ([xshift=1.1cm]0,0,3);
\end{tikzpicture}
\end{minipage}

\begin{bewijs}
Een punt $(x, y, z)$ voldoet aan $F(x, y) = 0$ als en slechts als zijn projectie $(x, y, 0)$ op de richtkromme ligt. Het punt ligt dan op de verticale rechte door een punt van de richtkromme.
\end{bewijs}
\end{theorie}

Zo stelt elke vergelijking zonder $x$ een cilinder voor met beschrijvenden evenwijdig met de $x$-as, en zonder $y$ een cilinder evenwijdig met de $y$-as. Een vlak $x + y = 2$ is ook zo'n cilinder: de richtkromme is een rechte. Dat zag je in Ruimtemeetkunde bij vlakken evenwijdig met een as.

\voorbeeld{}{Beschrijf de oppervlakken $z = y^2$ en $x^2 + z^2 = 1$.\label{vb:h2-cilinder}}{%
\uitleg{In de eerste vergelijking ontbreekt $x$, in de tweede $y$. Het zijn dus cilinders: we zoeken de richtkromme en de richting van de beschrijvenden.}
\naast{\begin{stappenvrij}{3.6cm}{5.4cm}
z = y^2 & Een parabool in het $yz$-vlak, met beschrijvenden evenwijdig met de $x$-as: een \textit{paraboolcilinder}, zoals een goot (rechts). \\[1ex]
x^2 + z^2 = 1 & Een cirkel in het $xz$-vlak, met beschrijvenden evenwijdig met de $y$-as: een buis langs de $y$-as.
\end{stappenvrij}}{\begin{tikzpicture}[zicht=.62,font=\scriptsize]
  \draw[->,lijngrijs] (0,0,0) -- (0,0,2.9) node[above,text=black] {$z$};
  \draw[->,lijngrijs] (0,-1.9,0) -- (0,2,0) node[right,text=black] {$y$};
  \fill[accent!15] plot[domain=-1.5:1.5,samples=30,smooth] (0,\x,{(\x)^2}) -- plot[domain=1.5:-1.5,samples=30,smooth] (3,\x,{(\x)^2}) -- cycle;
  \foreach \t in {-1.5,-1,...,1.5}{\draw[accent!60] (0,\t,{(\t)^2}) -- (3,\t,{(\t)^2});}
  \draw[accent,thick] plot[domain=-1.5:1.5,samples=30,smooth] (0,\x,{(\x)^2});
  \draw[accent,thick] plot[domain=-1.5:1.5,samples=30,smooth] (3,\x,{(\x)^2});
  \draw[->,lijngrijs] (3,0,0) -- (4,0,0) node[below left=-1pt,text=black] {$x$};
\end{tikzpicture}}}

\begin{oefeningen}
\oef{h2-cilinder}
\end{oefeningen}

\subsection{Omwentelingsoppervlakken}

Een pottenbakker draait een vaas op een draaischijf: elk punt van het profiel beschrijft een horizontale cirkel rond de as. Zo ontstaat een omwentelingsoppervlak. Welke vergelijking hoort daarbij? Neem een kromme in het $yz$-vlak met $y \geqslant 0$, en wentel ze om de $z$-as. Een punt $(0, y_0, z_0)$ van de kromme beschrijft de cirkel op hoogte $z_0$ met straal $y_0$: de punten met $z = z_0$ en $\sqrt{x^2 + y^2} = y_0$. De afstand tot de $z$-as neemt dus de plaats in van $y$.

\begin{theorie}
\begin{stelling}{Omwentelingsoppervlak}
Wentel je de kromme $F(y, z) = 0$ (met $y \geqslant 0$) in het $yz$-vlak om de $z$-as, dan krijg je het oppervlak

\vspace{1mm}
\qquad $\formule{F\bigl(\sqrt{x^2 + y^2},\, z\bigr) = 0}$
\end{stelling}

\begin{bewijs}
Een punt $P(x, y, z)$ ligt op de afstand $\rho = \sqrt{x^2 + y^2}$ van de $z$-as. Draai je $P$ om de $z$-as tot in het $yz$-vlak, aan de kant van $y \geqslant 0$, dan komt het in $(0, \rho, z)$. $P$ ligt op het omwentelingsoppervlak als en slechts als dat punt op de kromme ligt, dus als $F(\rho, z) = 0$.
\end{bewijs}
\end{theorie}

Je herkent een omwentelingsoppervlak om de $z$-as dus aan zijn vergelijking: $x$ en $y$ komen er alleen in voor via $x^2 + y^2$. Het oppervlak $z = x^2 + y^2$ van voorbeeld~\ref{vb:h2-doorsneden} is de parabool $z = y^2$, gewenteld om de $z$-as. Ook een bol is een omwentelingsoppervlak: wentel de halve cirkel $y^2 + z^2 = r^2$.

\voorbeeld{}{De \textit{sombrero} ontstaat als je de grafiek van $z = \dfrac{\sin y}{y}$ voor $y > 0$ om de $z$-as wentelt (met $z = 1$ op de as). Geef een vergelijking.\label{vb:h2-sombrero}}{%
\uitleg{Vervang $y$ door de afstand tot de as, $\sqrt{x^2 + y^2}$.}
\naast{\begin{stappenvrij}{4cm}{3.2cm}
z = \dfrac{\sin\sqrt{x^2 + y^2}}{\sqrt{x^2 + y^2}} & Voor $(x, y) \neq (0, 0)$.
\end{stappenvrij}
\vspace{1mm}
De golven van het profiel worden ringen. Elke nulwaarde $y = k\pi$ van het profiel geeft een cirkel met straal $k\pi$ in het $xy$-vlak. In de figuur is de hoogte vier keer uitvergroot.}{\begin{tikzpicture}[zicht=.22,font=\scriptsize]
  \beeld{sombrero.png}{-2.123}{-.759}{4.247}
\end{tikzpicture}}}

\voorbeeld{}{Een \textit{torus} (een donut) ontstaat als je de cirkel met middelpunt $(0, 2, 0)$ en straal 1 in het $yz$-vlak om de $z$-as wentelt. Geef een vergelijking.\label{vb:h2-torus}}{%
\uitleg{De cirkel is $(y - 2)^2 + z^2 = 1$. Hij ligt helemaal aan de kant $y > 0$, dus de stelling is van toepassing.}
\naast{\begin{stappenvrij}{4.6cm}{3.2cm}
\Bigl(\sqrt{x^2 + y^2} - 2\Bigr)^2 + z^2 = 1 & $y$ vervangen door de afstand tot de as.
\end{stappenvrij}
\vspace{1mm}
Controle: $(3, 0, 0)$ is het buitenste punt en $(1, 0, 0)$ het binnenste op de $x$-as; ze voldoen, want $(3 - 2)^2 = (1 - 2)^2 = 1$. Ook $(0, 2, 1)$, bovenaan, voldoet.}{\begin{tikzpicture}[zicht=.45,font=\scriptsize]
  \beeld{torus.png}{-1.400}{-.808}{2.800}
\end{tikzpicture}}}

\begin{werkwijze}
Een kromme in het $yz$-vlak wentelen om de $z$-as: vervang $y$ in haar vergelijking door $\sqrt{x^2 + y^2}$, de afstand tot de as. Omgekeerd: komen $x$ en $y$ alleen voor als $x^2 + y^2$, vervang dan $\sqrt{x^2 + y^2}$ door $y$ (met $y \geqslant 0$), en je vindt het profiel.
\end{werkwijze}

\nujij{Welk oppervlak krijg je als je de rechte $z = 2y$ ($y \geqslant 0$) om de $z$-as wentelt?}{$z = 2\sqrt{x^2 + y^2}$: een halve kegel met de top in de oorsprong.}

\begin{oefeningen}
\oef{h2-omwenteling}
\end{oefeningen}

%========================================================================================
\section{Parametervoorstellingen van oppervlakken}

\subsection{Twee parameters}

Met een vergelijking kun je nagaan of een punt op een oppervlak ligt, maar ze levert zelf geen punten. Een computer die een oppervlak tekent, heeft punten nodig. In Ruimtemeetkunde kreeg je elk punt van een vlak als $A + r\,\vec u + s\,\vec v$: twee getallen $r$ en $s$ geven één punt. Bij een gebogen oppervlak gaat dat ook, met andere formules.

\begin{theorie}
Een \textbf{\textcolor{accent}{parametervoorstelling}} van een oppervlak geeft de punten als
\bewijsregel{\formule{x = f(u, v), \quad y = g(u, v), \quad z = h(u, v)}}
met twee \textbf{\textcolor{accent}{parameters}} $u$ en $v$. Houd je $v$ vast en laat je $u$ variëren, dan krijg je een kromme op het oppervlak: een \textbf{\textcolor{accent}{parameterlijn}}. Zo ook met $u$ vast.
\end{theorie}

De parameterlijnen vormen samen een net op het oppervlak. De figuren in dit boek zijn zo getekend: de lijnen op de sombrero en de torus zijn parameterlijnen.

\voorbeeld{}{Geef een parametervoorstelling van de cilinder $x^2 + y^2 = 4$. Welke krommen zijn de parameterlijnen?\label{vb:h2-cilpar}}{%
\uitleg{Een cirkel met straal 2 heeft de parametervoorstelling $(2\cos u, 2\sin u)$. De hoogte $z$ is vrij: die wordt de tweede parameter.}
\begin{stappen}
x = 2\cos u, \quad y = 2\sin u, \quad z = v & Met $u \in [0, 2\pi[$ en $v \in \R$. \\[1ex]
x^2 + y^2 = 4\cos^2 u + 4\sin^2 u = 4 & Controle: elk punt ligt op de cilinder.
\end{stappen}
Met $v$ vast krijg je een horizontale cirkel op hoogte $v$. Met $u$ vast krijg je een verticale rechte: een beschrijvende.}

\subsection{De bol met lengte en breedte}

Op een wereldbol geef je een plaats met twee getallen: de breedte $\varphi$, ten noorden of ten zuiden van de evenaar, en de lengte $\lambda$, ten oosten of ten westen van de meridiaan van Greenwich. Dat zijn twee parameters. Welke coördinaten horen erbij? Het punt ligt op hoogte $z = R\sin\varphi$, op de breedtecirkel met straal $R\cos\varphi$, en op die cirkel bepaalt $\lambda$ de plaats.

\begin{theorie}
\begin{minipage}[t]{.64\linewidth}
\vspace{0pt}
\begin{stelling}{Parametervoorstelling van een bol}
De bol $x^2 + y^2 + z^2 = R^2$ heeft de parametervoorstelling

\vspace{1mm}
\quad $\formule{x = R\cos\varphi\cos\lambda, \ \ y = R\cos\varphi\sin\lambda, \ \ z = R\sin\varphi}$

\vspace{1mm}
met $\varphi \in \bigl[-\tfrac\pi2, \tfrac\pi2\bigr]$ de \textbf{\textcolor{accent}{breedte}} en $\lambda \in \left]-\pi, \pi\right]$ de \textbf{\textcolor{accent}{lengte}}.
\end{stelling}
\end{minipage}\hfill
\begin{minipage}[t]{.32\linewidth}
\vspace{0pt}
\centering
\begin{tikzpicture}[zicht=.9,font=\scriptsize]
  \bolzicht{.9}{0}{0}{0}{2}
  \draw[grijs,->] (0,0,0) -- (2.6,0,0) node[below left=-1pt,text=black] {$x$};
  \draw[grijs] (0,0,0) -- ({2*cos(40)},{2*sin(40)},0) -- ({2*cos(40)*cos(50)},{2*cos(50)*sin(40)},{2*sin(50)}) -- cycle;
  \draw[accent!70!black] plot[domain=0:40,samples=20,smooth,variable=\t] ({.6*cos(\t)},{.6*sin(\t)},0);
  \node at ({.75*cos(20)},{.75*sin(20)},0) {$\lambda$};
  \node at ({.55*cos(40)},{.55*sin(40)},.45) {$\varphi$};
  \fill[accent] ({2*cos(40)*cos(50)},{2*cos(50)*sin(40)},{2*sin(50)}) circle (1.6pt) node[above right=-1pt] {$P$};
\end{tikzpicture}
\end{minipage}

\begin{bewijs}
Elk zo'n punt ligt op de bol, want
\bewijsregel{x^2 + y^2 + z^2 = R^2\cos^2\varphi\,(\cos^2\lambda + \sin^2\lambda) + R^2\sin^2\varphi = R^2(\cos^2\varphi + \sin^2\varphi) = R^2 .}
Omgekeerd: neem een punt van de bol. Omdat $-R \leqslant z \leqslant R$, is er een $\varphi \in \bigl[-\tfrac\pi2, \tfrac\pi2\bigr]$ met $z = R\sin\varphi$. Dan is $x^2 + y^2 = R^2\cos^2\varphi$, dus $(x, y)$ ligt op de cirkel met straal $R\cos\varphi$, en is er een $\lambda$ met $x = R\cos\varphi\cos\lambda$ en $y = R\cos\varphi\sin\lambda$.
\end{bewijs}
\end{theorie}

De parameterlijnen zijn bekend: met $\varphi$ vast krijg je een breedtecirkel, met $\lambda$ vast een meridiaan, van pool tot pool. In de polen ($\varphi = \pm\tfrac\pi2$) hoort bij elke $\lambda$ hetzelfde punt: daar komen alle meridianen samen.

\voorbeeld{}{Brussel ligt op $50{,}85^\circ$ noorderbreedte en $4{,}35^\circ$ oosterlengte. Bereken de coördinaten, met $R = 6371$ km.\label{vb:h2-brussel}}{%
\uitleg{Een rechtstreekse invulling. Oosterlengte is positief, noorderbreedte ook.}
\begin{stappen}
x = 6371\cos 50{,}85^\circ\cos 4{,}35^\circ \approx 4011 & In km. \\[1ex]
y = 6371\cos 50{,}85^\circ\sin 4{,}35^\circ \approx 305 & Klein, want Brussel ligt dicht bij de meridiaan van Greenwich. \\[1ex]
z = 6371\sin 50{,}85^\circ \approx 4941 & Hoger dan $x$: Brussel ligt dichter bij de pool dan bij de evenaar.
\end{stappen}}

\begin{oefeningen}
\oef{h2-bol}
\end{oefeningen}

\subsection{Omwentelingsoppervlakken en de torus}

Dezelfde gedachte werkt voor elk omwentelingsoppervlak. Een punt van het profiel op afstand $\rho$ van de $z$-as en op hoogte $z$ beschrijft de cirkel $(\rho\cos u, \rho\sin u, z)$. Beschrijf je het profiel zelf met een parameter $v$, dan heb je twee parameters.

\begin{theorie}
\begin{stelling}{Parametervoorstelling van een omwentelingsoppervlak}
Wentel je de kromme $y = f(v)$, $z = g(v)$ (met $f(v) \geqslant 0$) in het $yz$-vlak om de $z$-as, dan krijg je

\vspace{1mm}
\qquad $\formule{x = f(v)\cos u, \quad y = f(v)\sin u, \quad z = g(v)} \qquad (u \in [0, 2\pi[)$
\end{stelling}
\end{theorie}

Voor $u = \tfrac\pi2$ krijg je $(0, f(v), g(v))$: het profiel zelf. De andere waarden van $u$ draaien het profiel rond. De bol met lengte en breedte is een geval: het profiel is de halve cirkel $y = R\cos\varphi$, $z = R\sin\varphi$.

\voorbeeld{}{Geef een parametervoorstelling van de torus uit voorbeeld~\ref{vb:h2-torus}, en controleer ze met zijn vergelijking.\label{vb:h2-toruspar}}{%
\uitleg{Het profiel is de cirkel met middelpunt $(0, 2, 0)$ en straal 1 in het $yz$-vlak: $y = 2 + \cos v$, $z = \sin v$.}
\begin{stappenbreed}
x = (2 + \cos v)\cos u, \ \ y = (2 + \cos v)\sin u, \ \ z = \sin v & Met $f(v) = 2 + \cos v \geqslant 1$. \\[1ex]
\sqrt{x^2 + y^2} = 2 + \cos v & De afstand tot de as. \\[1ex]
\bigl(\sqrt{x^2 + y^2} - 2\bigr)^2 + z^2 = \cos^2 v + \sin^2 v = 1 & De vergelijking klopt.
\end{stappenbreed}
De parameterlijnen met $u$ vast zijn cirkels rond de buis, die met $v$ vast horizontale cirkels rond de $z$-as.}

\begin{werkwijze}
Van een parametervoorstelling naar een cartesiaanse vergelijking: werk de parameters weg. Zoek combinaties zoals $x^2 + y^2$ waarin $\cos^2 + \sin^2 = 1$ een parameter doet verdwijnen, of druk een parameter uit in $x$, $y$ of $z$ en vul in.
\end{werkwijze}

\begin{oefeningen}
\oef{h2-param}
\end{oefeningen}

\subsection{De helicoïde}

Neem een horizontale rechte door de $z$-as, en laat ze gelijkmatig draaien terwijl ze gelijkmatig stijgt. Ze veegt een oppervlak uit dat je ziet in een wenteltrap, een schroef of de oprit van een parkeergarage: de \textit{helicoïde}. Is $v$ de draaihoek, dan staat de rechte op hoogte $cv$, en is $u$ de plaats op de rechte.

\begin{theorie}
\begin{minipage}[t]{.7\linewidth}
\vspace{0pt}
De \textbf{\textcolor{accent}{helicoïde}} is het oppervlak met parametervoorstelling

\vspace{1mm}
\qquad $\formule{x = u\cos v, \quad y = u\sin v, \quad z = cv} \qquad (u, v \in \R)$

\vspace{1mm}
met $c > 0$ een constante. Per volle draai stijgt ze $2\pi c$.
\end{minipage}\hfill
\begin{minipage}[t]{.26\linewidth}
\vspace{0pt}
\centering
\begin{tikzpicture}[zicht=.75,font=\scriptsize]
  \beeld{helicoide.png}{-1.550}{-.494}{3.100}
  \draw[->,lijngrijs] (0,0,3.14) -- (0,0,3.8) node[above,text=black] {$z$};
\end{tikzpicture}
\end{minipage}
\end{theorie}

De parameterlijnen met $v$ vast zijn de horizontale rechten. Die met $u$ vast zijn krommen die rond de $z$-as kronkelen op afstand $|u|$: schroeflijnen, het onderwerp van de volgende sectie.

\voorbeeld{}{Zoek een cartesiaanse vergelijking van de helicoïde met $c = 1$.\label{vb:h2-helicoide}}{%
\uitleg{Uit $z = v$ halen we $v$. Dan zeggen $x = u\cos z$ en $y = u\sin z$ dat $(x, y)$ op de rechte door de oorsprong met richting $(\cos z, \sin z)$ ligt. Dat wil zeggen dat $(x, y)$ en $(\cos z, \sin z)$ evenwijdig zijn.}
\begin{stappenbreed}
x\sin z - y\cos z = u\cos z\sin z - u\sin z\cos z = 0 & Controle met de parametervoorstelling. \\[1ex]
x\sin z = y\cos z & De vergelijking.
\end{stappenbreed}
Omgekeerd voldoet elk punt met $x\sin z = y\cos z$: dan zijn $(x, y)$ en $(\cos z, \sin z)$ evenwijdig, dus is $(x, y) = u(\cos z, \sin z)$ voor een getal $u$. Deze vergelijking is niet van de tweede graad: de helicoïde is geen kwadriek.}

%========================================================================================
\section{Ruimtekrommen}

\subsection{Een kromme met één parameter}

Een rechte in de ruimte heeft parametervergelijkingen met één parameter, of twee cartesiaanse vergelijkingen: ze is de snijlijn van twee vlakken. Zo is het ook met gebogen krommen in de ruimte. Eén parameter geeft een punt dat beweegt; twee oppervlakken snijden elkaar volgens een kromme.

\begin{theorie}
Een \textbf{\textcolor{accent}{ruimtekromme}} is een kromme met een parametervoorstelling
\bewijsregel{\formule{x = f(t), \quad y = g(t), \quad z = h(t)}}
met één parameter $t$. Een kromme die niet in één vlak ligt, heet \textbf{\textcolor{accent}{scheef}}.
\end{theorie}

Lees $t$ als de tijd, dan is de kromme de baan van een punt dat beweegt. Dat ken je van vlakke krommen, zoals de cirkel $(\cos t, \sin t)$; in de ruimte komt er gewoon een coördinaat bij.

\voorbeeld{}{Toon aan dat de kromme $x = 1 + \cos t$, $y = \sin t$, $z = 2\sin\tfrac t2$ (met $t \in [0, 4\pi]$) op de bol $x^2 + y^2 + z^2 = 4$ en op de cilinder $(x - 1)^2 + y^2 = 1$ ligt.\label{vb:h2-viviani}}{%
\uitleg{We vullen in, en gebruiken $\cos^2 + \sin^2 = 1$ en de formule $1 + \cos t = 2\cos^2\tfrac t2$.}
\naast{\begin{stappenvrij}{6.4cm}{2.5cm}
(x - 1)^2 + y^2 = \cos^2 t + \sin^2 t = 1 & Op de cilinder. \\[1ex]
x^2 + y^2 = 2 + 2\cos t = 4\cos^2\tfrac t2 & Uitgewerkt. \\[1ex]
x^2 + y^2 + z^2 = 4\cos^2\tfrac t2 + 4\sin^2\tfrac t2 = 4 & Op de bol.
\end{stappenvrij}}{\begin{tikzpicture}[zicht=.75,font=\scriptsize]
  \bolzicht[accent!45]{.75}{0}{0}{0}{2}
  \cirkelzicht[accent!45]{.75}{1}{0}{0}{1}
  \draw[accent,very thick] plot[domain=0:720,samples=120,smooth,variable=\t] ({1+cos(\t)},{sin(\t)},{2*sin(\t/2)});
\end{tikzpicture}}
Deze snijlijn van een bol en een cilinder heet de \textit{kromme van Viviani}. Ze ziet eruit als een acht die over de bol is gebogen.}

\begin{oefeningen}
\oef{h2-kromme}
\end{oefeningen}

\subsection{De kubische schroeflijn}

Een rechte en een vlak snijden elkaar hoogstens één keer, tenzij de rechte in het vlak ligt. Een kegelsnede en een rechte in haar vlak snijden elkaar hoogstens twee keer. De eenvoudigste scheve kromme is
\bewijsregel{x = t, \quad y = t^2, \quad z = t^3 ,}
de \textit{kubische schroeflijn}. Haar projectie op het $xy$-vlak is de parabool $y = x^2$, op het $xz$-vlak de kubische grafiek $z = x^3$. Hoe vaak kan een vlak haar snijden?

\begin{theorie}
\begin{stelling}{Vlak en kubische schroeflijn}
Een vlak snijdt de kubische schroeflijn $(t, t^2, t^3)$ in hoogstens drie punten. Vier verschillende punten van de kromme liggen dus nooit in één vlak.
\end{stelling}

\begin{bewijs}
Neem een vlak $ax + by + cz = d$, met $a$, $b$, $c$ niet alle nul. Het punt met parameter $t$ ligt erin als en slechts als
\bewijsregel{ct^3 + bt^2 + at - d = 0 .}
Dat is een veeltermvergelijking van graad hoogstens 3 die niet de nulveelterm is, want $a$, $b$ en $c$ zijn niet alle nul. Ze heeft hoogstens drie oplossingen. Lagen vier punten van de kromme in één vlak, dan had die vergelijking er vier.
\end{bewijs}
\end{theorie}

Daarom heet de kromme \textit{kubisch}: ze is van graad drie, zoals een rechte van graad één is en een kegelsnede van graad twee. Ze is zeker scheef, want een vlak dat haar bevat, zou oneindig veel punten met haar gemeen hebben.

\voorbeeld{}{Bepaal de snijpunten van de kubische schroeflijn met het vlak $11x - 6y + z = 6$.\label{vb:h2-kubisch}}{%
\uitleg{Vul $(t, t^2, t^3)$ in, en zoek de nulwaarden van de veelterm met Horner of door te proberen.}
\begin{stappen}
t^3 - 6t^2 + 11t - 6 = 0 & Ingevuld. \\[1ex]
(t - 1)(t - 2)(t - 3) = 0 & $t = 1$ is een nulwaarde; Horner geeft de rest. \\[1ex]
(1, 1, 1),\ (2, 4, 8),\ (3, 9, 27) & De drie snijpunten.
\end{stappen}}

\begin{oefeningen}
\oef{h2-kubisch}
\end{oefeningen}

\subsection{De schroeflijn}

De draad van een veer draait rond een as op een vaste afstand $R$, en stijgt gelijkmatig. Is $t$ de draaihoek, dan is de hoogte $ct$. Zo'n kromme ligt op de cilinder $x^2 + y^2 = R^2$, want $x$ en $y$ beschrijven een cirkel.

\begin{theorie}
De \textbf{\textcolor{accent}{schroeflijn}} (of \textbf{\textcolor{accent}{helix}}) met straal $R$ is de kromme
\bewijsregel{\formule{x = R\cos t, \quad y = R\sin t, \quad z = ct}}
met $c > 0$. De stijging per draai, $\formule{h = 2\pi c}$, heet de \textbf{\textcolor{accent}{spoed}}.

\begin{minipage}[t]{.6\linewidth}
\vspace{0pt}
\begin{stelling}{Schroeflijn uitgerold}
Rol je de cilinder uit tot een vlak, dan wordt de schroeflijn een rechte. Eén winding heeft dus de lengte

\vspace{1mm}
\qquad $\formule{L = \sqrt{(2\pi R)^2 + h^2}}$,

\vspace{1mm}
en de schroeflijn stijgt overal onder dezelfde hoek $\alpha$ met $\formule{\tan\alpha = \dfrac{h}{2\pi R}}$.
\end{stelling}
\end{minipage}\hfill
\begin{minipage}[t]{.36\linewidth}
\vspace{0pt}
\centering
\begin{tikzpicture}[zicht=.48,font=\scriptsize]
  \cirkelzicht[accent!45]{.48}{0}{0}{0}{1.3}
  \cirkelzicht[accent!45]{.48}{0}{0}{4.4}{1.3}
  \draw[accent!45] ([xshift=-.624cm]0,0,0) -- ([xshift=-.624cm]0,0,4.4) ([xshift=.624cm]0,0,0) -- ([xshift=.624cm]0,0,4.4);
  \draw[accent,thick] plot[domain=0:720,samples=100,smooth,variable=\t] ({1.3*cos(\t)},{1.3*sin(\t)},{2.2*\t/360});
\end{tikzpicture}\ \begin{tikzpicture}[scale=.36,font=\scriptsize]
  \draw[lijngrijs,fill=accent!8] (0,0) rectangle (8.17,4.4);
  \draw[lijngrijs] (0,2.2) -- (8.17,2.2);
  \draw[accent,thick] (0,0) -- (8.17,2.2) (0,2.2) -- (8.17,4.4);
  \node[below] at (4.08,0) {$2\pi R$};
  \node[right] at (8.17,1.1) {$h$};
  \draw[grijs] (1.6,0) arc[start angle=0,end angle=15.07,radius=1.6];
  \node at (2.2,.27) {$\alpha$};
\end{tikzpicture}
\end{minipage}

\begin{bewijs}
Snijd de cilinder open langs de verticale rechte door $(R, 0, 0)$ en rol hem uit. Een punt op de cilinder met draaihoek $t$ en hoogte $z$ komt in het vlak in $(Rt, z)$, want de boog van $0$ tot $t$ op de cirkel heeft lengte $Rt$. Het punt van de schroeflijn met parameter $t$ komt zo in $(Rt, ct)$: dat ligt op de rechte $z = \tfrac{c}{R}\,s$ in het uitgerolde vlak, met $s$ de horizontale coördinaat. Bij uitrollen veranderen lengtes op de cilinder niet. Eén winding wordt de schuine zijde van een rechthoekige driehoek met rechthoekszijden $2\pi R$ en $h$, en de helling van die rechte is $\tfrac{c}{R} = \tfrac{h}{2\pi R}$.
\end{bewijs}
\end{theorie}

Dat lengtes bij uitrollen niet veranderen, voel je aan met een vel papier: je kunt het oprollen tot een cilinder zonder het uit te rekken of te kreuken. Met een bol lukt dat niet: een kaart van de aarde vervormt altijd. Dat de lengte klopt, kun je ook narekenen met de booglengte uit Integralen (Koppel).

\voorbeeld{}{Terug naar het probleem. Een wenteltrap draait rond een zuil met een straal van 0,2~m. De treden zijn 1,2~m breed, dus de buitenrand ligt op 1,4~m van de as. Bij één volle draai stijgt de trap 3~m. Langs de buitenrand loopt een leuning. Hoe lang is de leuning voor één volle draai? En waarom is zo'n trap aan de binnenkant veel steiler dan aan de buitenkant: hoe steil is hij aan beide kanten, in graden?\label{vb:h2-terug}}{%
\uitleg{De leuning is een schroeflijn met $R = 1{,}4$ en spoed $h = 3$. De onderkant van de trap is een stuk van een helicoïde: elke trede ligt langs een horizontale rechte door de as. Op afstand $R$ van de as loopt over de treden een schroeflijn met dezelfde spoed, maar met een andere straal.}
\begin{stappenbreed}
L = \sqrt{(2\pi \cdot 1{,}4)^2 + 3^2} \approx \sqrt{77{,}4 + 9} \approx 9{,}29 & De leuning: ongeveer 9,3~m per draai. \\[2ex]
\tan\alpha = \dfrac{3}{2\pi \cdot 1{,}4} \approx 0{,}341, \quad \alpha \approx 18{,}8^\circ & Aan de buitenrand. \\[2ex]
\tan\alpha = \dfrac{3}{2\pi \cdot 0{,}2} \approx 2{,}39, \quad \alpha \approx 67{,}3^\circ & Aan de zuil.
\end{stappenbreed}
Aan de zuil leg je per draai maar $2\pi \cdot 0{,}2 \approx 1{,}26$~m af, en stijg je toch 3~m. Daarom loop je een wenteltrap aan de buitenkant.}

Zo zie je meteen een verschil tussen de helicoïde en een vlak. Op een hellend vlak is de helling overal dezelfde. Op de helicoïde hangt ze af van de afstand tot de as: het oppervlak is gedraaid.

\begin{werkwijze}
Een schroeflijn: zoek de straal $R$ en de spoed $h = 2\pi c$. Rol de cilinder uit: een winding wordt de schuine zijde van een rechthoekige driehoek met rechthoekszijden $2\pi R$ en $h$.
\end{werkwijze}

\begin{oefeningen}
\oef{h2-schroef}
\end{oefeningen}

%========================================================================================
\opdrachtenkop

\begin{oefeningen}
\oef{h2-o1}

\oef{h2-o2}

\oef{h2-o3}

\oef{h2-o4}

\oef{h2-o5}

\oef{h2-o6}

\oef{h2-o7}

\oef{h2-o8}

\oef{h2-o9}

\oef{h2-o10}

\oef{h2-o11}

\oef{h2-o12}

\oef{h2-o13}

\oef{h2-o14}

\oef{h2-o15}

\oef{h2-o16}

\oef{h2-o17}

\oef{h2-o18}

\oef{h2-o19}

\oef{h2-o20}

\oef{h2-o21}

\oef{h2-o22}

\oef{h2-o23}
\end{oefeningen}

%----------------------------------------------------------------------------------------
% KOPPEL: verbanden met andere leerstof
%----------------------------------------------------------------------------------------
\needspace{20\baselineskip}
\opdrachtsectie{Koppel}

\lkr{Integralen staat in het 6de jaar: geef de eerste en de tweede oefening pas na het hoofdstuk over toepassingen (booglengte en volume van een omwentelingslichaam). Complexe getallen zag je in het 5de jaar.}

\begin{oefeningen}
\oef{h2-koppel-1}

\oef{h2-koppel-2}

\oef{h2-koppel-3}
\end{oefeningen}

%----------------------------------------------------------------------------------------
\needspace{30\baselineskip} % de test en het schema Fout? Herhaal samen op een pagina
\opdrachtsectie{Test jezelf}

\begin{oefeningen}
\oef{h2-t1}

\oef{h2-t2}

\oef{h2-t3}

\oef{h2-t4}

\oef{h2-t5}

\oef{h2-t6}

\oef{h2-t7}
\end{oefeningen}

\herhaalschema{%
  \fv{oef:h2-t1}{vb:h2-doorsneden}{oef:h2-doorsneden}%
  \fv{oef:h2-t2}{vb:h2-cilinder}{oef:h2-cilinder}%
  \fv{oef:h2-t3}{vb:h2-torus}{oef:h2-omwenteling}%
  \fv{oef:h2-t4}{vb:h2-toruspar}{oef:h2-param}%
  \fv{oef:h2-t5}{vb:h2-brussel}{oef:h2-bol}%
  \fv{oef:h2-t6}{vb:h2-kubisch}{oef:h2-kubisch}%
  \fv{oef:h2-t7}{vb:h2-terug}{oef:h2-schroef}}

\ifzelfstudie
\section*{Maak je eigen samenvatting}
Sluit dit boek en schrijf in je schrift: wat een cartesiaanse vergelijking en een niveaukromme zijn, waarom een vergelijking zonder $z$ een cilinder is, hoe je een omwentelingsoppervlak herkent en opstelt, de parametervoorstellingen van de cilinder, de bol, een omwentelingsoppervlak en de helicoïde, wat een ruimtekromme is, waarom een vlak de kubische schroeflijn hoogstens drie keer snijdt, en de lengte en helling van een schroeflijn. Vergelijk daarna met de kaders.
\fi
